% H3 (amendment A1), source: results/phase2/h3/analysis/final/{table_h3.md,h3_paired.csv}
\textbf{H3, battery-aware fallback (pre-registered amendment A1).} H3 was not covered by the original protocol. We added
it after the main campaign as a pre-registered amendment, committed and frozen before any H3 data existed. The amended
code reproduced a frozen v1.1 cell byte-for-byte. The amendment compares the fallback with an \emph{energy-blind}
variant. This variant uses the same planner with unit cell costs, a constant unknown-cell premium and no route
feasibility check, while the monitor, gates, thresholds and maps are identical. The pre-registered primary metric is
the share of fallback episodes that do not end in battery depletion. It rises significantly only under combined shift,
and only by 0.2 points (Holm $p=0.039$), so H3 is formally supported. This metric is near its ceiling, however, because
safe stop prevents depletion with either planner. The effect appears in the secondary metrics (Fig.~\ref{fig:h3}). With
the energy-blind fallback, success under combined shift collapses from 0.71 to 0.27 ($+44.5$ points for battery-aware;
73\% of episodes end in a stop once the monitor vetoes moves the battery can no longer afford). Energy use rises by 33--40
units under sensor and combined shift, where the fallback is in control most of the time. Without safe stop, the
battery-aware fallback avoids depletion in 47 points more of the fallback episodes under combined shift. In-distribution
the energy-blind planner is marginally better (success $+1.0$ point), because it stops less often. Energy awareness in
the fallback therefore matters when the fallback carries the mission under harsh shift, and costs a little when it
rarely acts.

\begin{figure}[t]
\centering
\includegraphics[width=\columnwidth]{fig8_h3}
\caption{Amendment A1 (H3): paired differences, battery-aware $-$ energy-blind fallback, with and without safe stop
(mean, 95\% CI over 10 seeds; * Holm $p<0.05$).}
\label{fig:h3}
\end{figure}
